\bookchapter{Kapitel 32: Notfallhandbuch für Smartphone-Sicherheit}

\chapterintro{Dieses Kapitel ist ein „Notfallplan“. Hoffentlich brauchen Sie ihn nie – falls doch, hilft er Ihnen, bei Verlust, Diebstahl oder Betrug den Schaden zu begrenzen. Auch bei Erdbeben, Starkregen und anderen Katastrophen ist das Smartphone Werkzeug für Warnungen, Standortnotruf und Kommunikation. Lesen Sie das Kapitel am besten jetzt einmal durch, speichern Sie die wichtigsten Schritte als Screenshot oder drucken Sie sie aus und legen Sie sie zu Hause hin – im Ernstfall ist das Handy womöglich weg, bei Katastrophen leer oder ohne Netz.}

\booksection{Szenario 0: Erdbeben / Extremwetter – wie das Smartphone Leben rettet}

Hier geht es um körperliche Sicherheit, nicht um Konten und Geld. Zu Warn-Einstellungen siehe Kapitel 13 „Wetterwarnungen, Naturgefahren und Erdbebenwarnung“; hier nur: Was tun, wenn der Alarm ertönt?

\textbf{Erdbebenalarm: Erst Schutz suchen, dann evakuieren – nicht nur filmen}\par

\begin{enumerate}
  \item In der Nähe Schutz suchen: Hinlegen, abdecken, festhalten (unter Tischen, an tragenden Wänden, an stabilen Möbeln), Kopf und Nacken schützen; von Glas, Außenwänden und Kronleuchtern weg.
  \item Nicht blind ins Treppenhaus stürmen: Das Hauptbeben dauert oft nur Sekunden bis Dutzende Sekunden, Treppenhäuser können gefährlicher sein. Aufzüge niemals benutzen.
  \item Nach dem Beben geordnet evakuieren: Treppe nutzen, auf Nachbeben und herabfallende Teile achten; auf offene Flächen, fern von Masten, Reklame und maroden Gebäuden.
  \item Handy gezielt nutzen: Kurze Nachrichten an Notfallkontakte („Bin bei XX, sicher/verletzt, brauche Hilfe“) sparen mehr Akku als lange Telefonate. Wenn SMS/WeChat möglich, lieber Text.
  \item Ohne Netz: siehe Kapitel 7 – Satelliten-SOS / Satellitennachrichten (falls das Gerät es unterstützt). Oder auf Anhöhe/offene Fläche gehen und erneut versuchen. In manchen Fällen hilft Bluetooth-Sprechfunk ohne Netz.
  \item Nicht lebenswichtige Funktionen abschalten: Im Notfall Flugmodus mit Zeitplan ein-/ausschalten, Bluetooth-Scans stoppen, Helligkeit senken – das bringt oft Stunden. Den letzten Akku für „Ich bin sicher“ und Rettungsnachrichten aufsparen.
\end{enumerate}

\begin{tcbwarn}{Starkregen / Taifun / Sturzfluten: Schon in der Warnphase handeln}
\begin{itemize}
  \item Bei roter Warnung: Möglichst drinnen bleiben, Unterführungen, Flussauen und Hänge meiden.
  \item Draußen eingeschlossen: In massive Gebäude höher hinaufgehen, nicht durch tiefes Wasser waten (fehlende Deckel, Stromschlaggefahr). Beim Notruf 119/110 Ort und Wasserstand klar nennen.
  \item Vor der Abreise: Risiko über Wetter-App und lokale Warnkanäle prüfen; bei langen Autofahrten Passagen an Hängen und Gewässern meiden.
  \item In der Familie vereinbaren: „Bei Kontaktverlust gilt der letzte Standort aus WeChat/SMS“ – so verhindert man endloses Wählen und leeren Akku.
\end{itemize}

\end{tcbwarn}

\begin{tcbpractice}{Vier Dinge, die Sie heute erledigen sollten}
\begin{enumerate}
  \item System-Erdbebenwarnung bzw. eine seriöse Warn-App einschalten und erlauben, den Stumm-Modus zu durchbrechen.
  \item 2–3 „Notfallkontakte“ in den Kontakten anheften, damit sie auf dem Sperrbildschirm erscheinen (iOS/Android: Medizinische Notfallkarte oder Sperrbildschirm-Infos).
  \item Powerbank zu Hause und in der Tasche; vor langen Reisen laden und Kabel mitnehmen.
  \item Mindestens eine Nummer auswendig können, die „ohne Kontaktliste“ erreichbar ist (Handy der Familie), auf Papier in die Geldbörse legen.
\end{enumerate}

\end{tcbpractice}

\booksection{Szenario 1: Handy verloren/gestohlen – die goldenen 30 Minuten}

Die ersten 30 Minuten nach Verlust oder Diebstahl sind die „goldene Zeit“. Richtiges Handeln schützt Geld und Privatsphäre. Reihenfolge einhalten, keinen Schritt auslassen:

\textbf{Standardablauf in den ersten 30 Minuten}\par

\begin{enumerate}
  \item Mit einem fremden Handy die eigene Nummer anrufen (Minute 1): Vielleicht ist es nur in der Nähe gefallen und jemand hebt ab.
  \item Fernsperre setzen (Minute 2–5): Über ein anderes Gerät den Suchdienst nutzen: iPhone: iCloud.com > Suchen > Gerät wählen > Als verloren markieren
  \item Android je Marke: Huawei cloud.huawei.com / Xiaomi i.mi.com / OPPO cloud.oppo.com usw.
  \item SIM sperren lassen (Minute 5–10): Kundendienst anrufen (China Mobile 10086 / Unicom 10010 / Telecom 10000) und SIM sperren lassen. Der wichtigste Schritt – wer Ihre SIM hat, empfängt Bestätigungscodes und kann alle Passwörter zurücksetzen.
  \item Zahlungskonten einfrieren (Minute 10–20): Alipay: 95188 anrufen oder auf einem fremden Gerät einloggen > Meine Seite > Einstellungen > Sicherheitscenter > Notfall-Sperre
  \item WeChat: Nach 110 den WeChat-Support kontaktieren oder im Sicherheitscenter > Konto sperren
  \item Bank-App: Kundendienst auf der Kartenrückseite anrufen und Online-Banking sperren lassen
  \item Angehörige informieren (Minute 20–25): Per Moments oder über Familie Nachrichten senden, dass das Handy weg ist – damit niemand in Ihrem Namen betrogen wird.
  \item Anzeige erstatten (Minute 25–30): 110 anrufen oder auf die nächste Wache gehen und eine Bestätigung erhalten. Für Versicherung und Bankstreitigkeiten brauchen Sie dieses Dokument.
\end{enumerate}

\booksection{Szenario 2: Betrogen worden – die goldenen 72 Stunden}

Nach einem Betrug: Ruhe bewahren. Panische Fehler führen zu weiteren Verlusten. So vorgehen:

\begin{enumerate}
  \item Sofort keine weiteren Überweisungen: Egal was gesagt wird – keinen Cent mehr. „Noch eine Sicherheitszahlung, dann kommt alles zurück“ ist der Klassiker des zweiten Betrugs.
  \item Alle Beweise sichern: Chatverläufe nicht löschen! Anrufprotokolle nicht löschen! Chats, Überweisungen und Kontodaten der Gegenseite screenshotten. Das ist entscheidend für Anzeige und Rückholung.
  \item Sofort 96110 anrufen: Die nationale Anti-Betrug-Hotline. Die Leitstelle hilft bei der Notfall-Sperre – wenn die Überweisung noch nicht lange her ist (meist innerhalb von 24 Stunden), kann die Bank das Gegenkonto sperren.
  \item 110 anrufen: Mit Handy, Ausweis und Überweisungsbeleg zur nächsten Wache und Protokoll aufnehmen lassen.
  \item Bank/Zahlungsdienst kontaktieren: Betrug melden und Kontosperre verlangen. Bei Kreditkartenmissbrauch die Bank um Chargeback bitten.
  \item Alle Passwörter ändern: Auf einem sicheren Gerät Apple-ID/Markenkonto, WeChat, Alipay und Bank-App-Passwörter erneuern und Zwei-Faktor aktivieren.
\end{enumerate}

\booksection{Szenario 3: Virus/Trojaner auf dem Handy – was tun?}

Warnzeichen: Plötzliche Werbung, automatisch installierte Apps, explodierende Telefonrechnung, ungewöhnlicher Datenverbrauch – vermutlich Schadsoftware:

\begin{itemize}
  \item Schritt 1: Abgesicherter Modus – Android: Ein-/Ausschalter lange drücken und ausschalten, dann wieder einschalten, nach dem Vibrieren Leiser-Taste gedrückt halten, bis „Abgesicherter Modus“ erscheint (unten links). Dort sind Drittanbieter-Apps deaktiviert – ideal zur Fehlersuche.
  \item Schritt 2: Verdächtige Apps deinstallieren – Im abgesicherten Modus zu Einstellungen > Apps, nach Installationsdatum sortieren und unbekannte kürzlich installierte Apps entfernen.
  \item Schritt 3: Sicherheits-App scannen – „Tencent Mobile Manager“ oder „360 Mobile Guard“ installieren und einen vollständigen Scan durchführen.
  \item Schritt 4: Werksreset (letzte Lösung) – Wenn nichts hilft: Wichtige Daten sichern und zurücksetzen. Achtung: Es werden alle Daten gelöscht – Fotos, Kontakte usw. vorher sichern.
  \item iPhone wird praktisch nie „krank“ – die iOS-Sandbox macht Schadsoftware schwer. Bei Anomalien steckt meist ein App-Bug oder eine Einstellung dahinter, kein Virus. Ein Reset löst 99 \% der Probleme.
\end{itemize}

\booksection{Szenario 4: Konto gekapert – Kontrolle zurückholen}

\begin{enumerate}
  \item Sofort Passwort ändern: Wenn Sie noch einloggen können, als erstes ein neues, völlig anderes Passwort vergeben.
  \item Andere Geräte entfernen: In den Sicherheitseinstellungen „Angemeldete Geräte“ suchen und unbekannte Geräte abmelden.
  \item Anmeldungen prüfen: Letzte 7 Tage kontrollieren. Anmeldungen von fremden Orten deuten auf Diebstahl hin.
  \item Zwei-Faktor einschalten: Falls noch nicht geschehen, sofort aktivieren – der wirksamste Schutz gegen erneuten Diebstahl.
  \item Verknüpfte Daten prüfen: Handynummer und E-Mail des Kontos kontrollieren. Bei Änderung über die Kontowiederherstellung zurückholen.
  \item Zahlungseinstellungen prüfen: Bei WeChat/Alipay ungeschützte Zahlungen und Lastschriften kontrollieren und Bankkarten auf unbekannte Konten prüfen.
\end{enumerate}

\booksection{Szenario 5: Erpressungsnachricht – keine Angst, nicht zahlen}

Zwei typische Maschen: (1) „Ich habe kompromittierende Videos/Fotos, sonst geht es an alle deine Kontakte“; (2) „Dein Handy ist gehackt, sonst wird es gesperrt“. Dann so handeln:

\begin{itemize}
  \item 99,9 \% sind Betrug: Es existieren keine sensiblen Daten. Dieselbe Nachricht geht an Zehntausende – irgendwer zahlt aus Angst.
  \item Nicht antworten, nicht zahlen: Eine Antwort zeigt „Diese Nummer wird benutzt“ und Sie landen auf der „erpressbaren“ Liste – danach wird es schlimmer.
  \item Screenshot als Beweis: Danach blockieren.
  \item Passwort ändern: Zur Sicherheit die wichtigsten Konten-Passwörter erneuern.
  \item Wenn echte Datenpanne droht: Anzeige erstatten. Die Polizei hat Erfahrung mit solchen Fällen.
\end{itemize}

\booksection{Szenario 6: KI-Gesichts-/Stimmklau – Augenzeugen sind nicht mehr Beweis}

Seit 2024 können Betrüger per KI-Gesicht und KI-Stimme im Videoanruf „zu“ Familie, Vorgesetzten oder Freunden werden. „Mama, mir ist was passiert, schick Geld“ – im Video sieht und klingt alles nach Ihrem Sohn, ist aber Fake. Eine der derzeit schwersten und gefährlichsten Methoden.

\textbf{Schritt 1: Schwachstellen erkennen – vier typische Mängel von KI-Fälschungen}\par

\begin{itemize}
  \item Ungewöhnliche Blinzelfrequenz: Menschen blinzeln 15–20 Mal pro Minute; KI-Gesichter blinzeln zu selten oder steif.
  \item Mund und Ton nicht synchron: Zwischen Lippenbewegung und Ton liegen oft Bruchteile einer Sekunde.
  \item Unscharfe Gesichtskonturen: An Wangenknochen, Ohren und Haaransatz treten „Ränder“ oder Flimmern auf.
  \item Licht inkonsistent: Glanzlichter und Schatten im Gesicht passen nicht zum Hintergrund, oder bei Seitenansicht „drehen“ sich Gesichtszüge falsch mit.
\end{itemize}

\begin{tcbinfo}{Schritt 2: Drei Prüfungen, ob die Person echt ist}
\begin{itemize}
  \item Bestimmte Bewegungen verlangen: Kopf drehen, Gesicht verdecken, Hand vor den Mund und dann Mund öffnen/schließen. Bei KI-Klau sind Verdeckungen und große Drehungen die Schwachstellen – das Gesicht verzerrt oder hinkt hinterher.
  \item Private Frage stellen: Eine erfundene, z. B. „In welche Tierklinik war unser Hund letztes Mal?“ – Betrüger stocken oder antworten falsch.
  \item Bekannte Nummer zurückrufen: Auflegen und die gespeicherte Nummer direkt anrufen. Echte Handynummern lassen sich nicht so leicht kapern – der Betrug fliegt auf.
\end{itemize}

\end{tcbinfo}

\begin{tcbwarn}{Schritt 3: Nach bestätigtem Betrug die Anzeige}
\begin{enumerate}
  \item Sofort keine Überweisung: Auch wenn es „dringend“ klingt – erst stoppen. Wer zu „schnell überweisen, nicht auflegen“ drängt, ist fast immer Betrüger.
  \item Alle Beweise sichern: Bildschirmaufnahme, Videochat, Chats, Konten, Überweisungsbelege – das ist entscheidend für die technische Begutachtung.
  \item 96110 anrufen: Als „KI-Gesicht/Stimme“ melden – die Leitstelle startet die Notfall-Sperre.
  \item 110 anrufen: Mit allen Beweisen zur Wache, ausdrücklich KI-Fälschung erwähnen, damit die Polizei den Fall richtig einordnen kann.
  \item Angehörige informieren: Betrüger könnten Ihr gefälschtes Video für weitere Taten nutzen – rechtzeitig in Moments oder Familienchats warnen.
\end{enumerate}

\end{tcbwarn}

\booksection{Szenario 7: SIM-Swap – Ihre Nummer wird „gestohlen“}

Eine wenig bekannte, aber extrem gefährliche Attacke. Beim SIM-Swap geben sich Angreifer mit gefälschten Identitätsdaten als Sie aus und lassen im Shop eine neue SIM ausstellen. Danach verliert Ihr Handy plötzlich den Empfang, während das fremde Gerät Ihre SMS-Codes empfängt – Alipay, Bank, WeChat stehen weit offen.

\begin{tcbinfo}{So erkennen Sie es: Drei Warnsignale}
\begin{itemize}
  \item Plötzlich „Kein Dienst“/„Keine SIM-Karte“, obwohl Sie die Karte nicht angetastet haben und ein Neustart nicht hilft – oft der Moment der Übernahme.
  \item SMS des Anbieters über „Kartenersatz erfolgreich“ oder „Dienst erfolgreich gebucht“ – obwohl Sie nichts veranlasst haben.
  \item Plötzlich keine Codes mehr, aber auch sonst keine SMS – weil die Nachrichten auf die fremde Karte umgeleitet werden.
\end{itemize}

\end{tcbinfo}

\begin{tcbwarn}{Schutzmaßnahmen – die SIM-Karte „abschließen“}
\begin{itemize}
  \item SIM-PIN setzen: Einstellungen > Telefon > SIM-PIN. Danach ist der PIN nach jedem Neustart oder Kartenwechsel nötig – auch wer die Karte hat, kann sie so nicht nutzen.
  \item Zweite Verifikation beim Anbieter: In der China Mobile/Unicom/Telecom-App „Sicherheitscenter“ oder „Dienste-Verifikation“ öffnen und Warnungen bei fremdem Kartenersatz sowie zweite Passwortprüfung aktivieren. So verlangt der Shop beim Ersatz eine Zusatzprüfung.
  \item Wichtige Konten nicht nur per SMS absichern: Für Bank, Apple ID usw. Authenticator-Apps (Google Authenticator, Microsoft Authenticator) oder Hardware-Schlüssel statt SMS-Codes verwenden.
  \item Ausweiskopien mit Wasserzeichen: Auf Kopien „Nur für XX-Zweck“ vermerken – Angreifer missbrauchen oft geleakte Ausweisdaten zur Identitätsfälschung.
\end{itemize}

\end{tcbwarn}

\textbf{Nach der Übernahme: Tempo ist alles}\par

\begin{enumerate}
  \item Sofort Anbieter anrufen und sperren lassen: Mobile 10086 / Unicom 10010 / Telecom 10000, sofortige Sperre und Stopp des Kartenersatzes verlangen.
  \item Auf anderen Geräten einloggen: Am PC oder Handy der Familie sofort Passwörter für Bank, Alipay, WeChat und Apple ID ändern und unbekannte Geräte abmelden.
  \item Zahlungskonten sperren: Alipay 95188, WeChat-Sicherheitscenter, Bank-Online-Banking sperren.
  \item Mit Ausweis zum Shop: Persönlich in einer offiziellen Filiale eine neue Karte ausstellen lassen und die Nummer zurückholen.
  \item Anzeige: SIM-Swap kann als Verstoß gegen den Schutz personenbezogener Daten gewertet werden. Mit Belegen (Sperrezeit, Schäden, Ersatzaufzeichnungen) anzeigen.
  \item Alles durchleuchten: Nach Wiederherstellung Anmeldungen und Geldbewegungen prüfen, ob Identitätsmissbrauch stattgefunden hat.
\end{enumerate}

\begin{tcbwarn}{Vorbeugen – fünf Dinge im Alltag}
\begin{enumerate}
  \item Bildschirmsperre + Zwei-Faktor, beides aktivieren
  \item „Handy suchen“ muss eingeschaltet sein
  \item SIM-PIN setzen (Einstellungen > Telefon > SIM-PIN) – bei Verlust kann niemand Ihre Karte nutzen
  \item Keine Apps unbekannter Herkunft, keine fremden Links
  \item Wichtige Daten regelmäßig sichern (Fotos, Kontakte, Chats)
\end{enumerate}

\end{tcbwarn}

\begin{tcbpractice}{Übung: Ihre persönliche Notfallkarte}
\begin{enumerate}
  \item Ein Blatt Papier nehmen und notieren: Anbieter-Support (Mobile 10086 / Unicom 10010 / Telecom 10000), Bank-Support, 96110 Anti-Betrug, 110, Telefonnummern der Familie.
  \item Registrierte E-Mail und Ersatzkontakt für Apple ID / Markenkonto notieren.
  \item Blatt in die Geldbörse oder an den Kühlschrank – nicht nur aufs Handy, denn im Notfall ist es womöglich weg.
  \item Jetzt prüfen: Ist „Handy suchen“ aktiviert, ist die SIM-PIN gesetzt?
\end{enumerate}

\end{tcbpractice}

\begin{tcbsummary}{Kapitel-Zusammenfassung}
\begin{itemize}
  \item Handy weg – goldene 30 Minuten: Fremdes Handy nutzen Fernsperre SIM sperren Zahlungen einfrieren Familie informieren Anzeige
  \item Betrug – 72 Stunden: Überweisungen stoppen Beweise sichern 96110 110 Bank Passwörter ändern
  \item Schadsoftware: Abgesicherter Modus Apps deinstallieren Scan Werksreset (letzte Lösung)
  \item Konto gekapert: Passwort Geräte abmelden Anmeldungen prüfen 2FA Bindungen prüfen Zahlungen prüfen
  \item Erpressung: 99,9 \% Betrug, nicht antworten/zahlen, screenshoten und blockieren
  \item KI-Gesicht/Stimme: Auf Blinzeln/Mund/Kanten/Licht achten, Kopfdrehung und Verdeckung verlangen, private Frage stellen, bekannte Nummer zurückrufen
  \item SIM-Swap: „Kein Dienst“ ist Warnsignal; SIM-PIN setzen, Anbieter-2FA, wichtige Konten mit Authenticator
  \item Fünf Präventionsmaßnahmen: Sperre+2FA, Handy suchen, SIM-PIN, keine unbekannten Apps, Backups
\end{itemize}

\end{tcbsummary}
